% -----------------------------------------------------------------------------
\section{Simulation Environment, Experiment Workflow, and Outputs}
\label{sec:environment}
% -----------------------------------------------------------------------------

MoonAI is centered on the simulation environment itself. The final runtime defines how
experiments are described, how they are executed, what data is produced, and how that data is
reviewed after the run.

\subsection{Experiment Definition and Parameter Space}

The base experiment structure is defined by \texttt{SimulationConfig}. The default
configuration includes a square world, separate predator and prey populations, food entities,
energy rules, mutation rates, compatibility parameters, reporting cadence, and seed control.
Current defaults include a world size of 3000, 12000 predators, 48000 prey, 60000 food items,
and a report interval of 1000 ticks.

The shipped \texttt{experiments.lua} file expands that base into a broad experiment matrix. It
defines 55 named conditions and applies 5 fixed seeds to each, producing 275 seeded runs, plus
an additional \texttt{default} entry for ad hoc use. The experiment groups cover multiple kinds
of study rather than only one benchmark dimension.

The main experiment groups are:

\begin{description}[style=nextline,leftmargin=4.1cm,font=\normalfont\bfseries]
  \item[Baseline sweeps] Vary mutation, speciation, and population size to compare broad
  evolutionary behavior under different default-scale pressures.

  \item[Scale experiments] Vary population and world size to observe how runtime behavior and
  population dynamics change with density and total count.

  \item[Density experiments] Hold population constant while changing world size to study sparse
  and dense ecological pressure.

  \item[Run-length experiments] Change \texttt{max\_ticks} to compare short and long runs under
  otherwise similar conditions.

  \item[Energy and resource experiments] Alter food respawn and energy budget parameters to stress
  survival and reproduction under scarce or abundant resources.

  \item[Perception and mobility experiments] Vary vision, speed, and interaction range to observe
  how sensing and movement shape survival strategies.

  \item[Complexity experiments] Vary add-node, add-connection, and hidden-node limits to study how
  network growth constraints affect behavior and topology.
\end{description}

\subsection{Interactive Runtime Workflow}

The application always launches the UI. From the interface, a user can select a preset, edit a
draft configuration, queue runs, replace the current run, inspect the world, and adjust visual
settings without modifying source files.

\begin{figure}[H]
  \centering
  \includegraphics[width=\textwidth,height=0.68\textheight,keepaspectratio]{app-screen.png}
  \caption{MoonAI desktop application showing the live simulation interface and experiment
  inspection workflow.}
  \label{fig:app-screen}
\end{figure}

The runtime view shown in Figure~\ref{fig:app-screen} reflects several design decisions:

\begin{itemize}[nosep]
  \item the full active population is rendered instead of a sampled subset,
  \item a selected agent can expose additional information such as sensor lines and neural
    topology,
  \item queue management is built into the application shell,
  \item settings are persisted separately from experiment definitions.
\end{itemize}

This separation between experiment definitions and UI state is important for repeatability.
Simulation parameters live in \texttt{experiments.lua}, while interface settings live in
\texttt{settings.json}.

\subsection{Selected-Agent Inspection and Observability}

MoonAI does more than display moving points. When an agent is selected, the UI can show a
vision circle, sensor lines, a stats panel, and a network panel. The controller receives 35
inputs consisting of nearest predators, nearest prey, nearest food items, self state, and wall
proximity. This makes the runtime useful not only for entertainment or visualization, but for
actual inspection of how evolved behavior is being produced.

The runtime also includes a scoped profiler tree. This profiler is useful when evaluating where
time is spent across spatial grid construction, movement, reproduction, and other host-side
control scopes.

\subsection{Output Artifacts and Analysis Pipeline}

Every run writes a structured output bundle under \texttt{output/experiments/}. The primary
artifacts are:

\begin{description}[style=nextline,leftmargin=3.5cm,font=\normalfont\ttfamily]
  \item[config.json] Full configuration snapshot for the executed run.
  \item[stats.csv] Time-series population and complexity metrics at each report interval.
  \item[species.csv] Species-level population summaries over time.
  \item[genomes.json] Representative genome snapshots including nodes and connections.
\end{description}

Those artifacts feed the Python analysis package. The analysis pipeline discovers valid runs,
groups them by condition, builds comparison plots, renders representative genomes, and writes a
self-contained HTML report.

\begin{figure}[H]
  \centering
  \includegraphics[width=\textwidth,height=0.68\textheight,keepaspectratio]{analysis-report.png}
  \caption{Generated HTML analysis report used to compare experiment conditions and summarize
  completed runs.}
  \label{fig:analysis-report}
\end{figure}

The report shown in Figure~\ref{fig:analysis-report} is an important part of the final system.
It turns raw exports into a usable comparison surface for researchers, which reinforces the
project's main goal: providing an environment for studying evolutionary machine learning, not
merely running one simulation loop.
